\documentclass[11pt,a4paper]{article}
\usepackage[T1]{fontenc}
\usepackage{lmodern}
\usepackage[margin=21mm]{geometry}
\usepackage{xcolor}
\usepackage{enumitem}
\usepackage{fancyhdr}
\definecolor{PrimaryBlue}{RGB}{0,0,180}
\definecolor{PrimaryRed}{RGB}{185,0,0}
\usepackage[colorlinks=true,urlcolor=PrimaryBlue,pdftitle={Curriculum Vitae - Soumya Swayamjyoti},pdfauthor={Soumya Swayamjyoti}]{hyperref}
\setlength{\parindent}{0pt}
\setlength{\parskip}{4pt}
\setlist{nosep,leftmargin=1.45em}
\setlength{\headheight}{14pt}
\pagestyle{fancy}
\fancyhf{}
\fancyfoot[L]{\small Soumya Swayamjyoti}
\fancyfoot[R]{\small\textcolor{PrimaryBlue}{\thepage}}
\renewcommand{\headrulewidth}{0pt}
\newcommand{\cvsection}[1]{\par\addvspace{9pt}{\large\sffamily\bfseries\color{PrimaryBlue}#1}\par\nobreak\vspace{1pt}}
\newcommand{\entry}[2]{\textbf{#1}\hfill\textcolor{PrimaryBlue}{#2}\par\nobreak}
\newcommand{\scholar}{\href{https://scholar.google.com/citations?user=w6ZvEUwAAAAJ\&hl=en}{Google Scholar}}

\begin{document}

aggedright
{\fontsize{30}{34}\selectfont\sffamily\bfseries\color{PrimaryBlue} Soumya Swayamjyoti}\par
{\large Computational Materials Scientist and Machine Learning Researcher}\par
Visiting Researcher, RWTH Aachen University \quad\textcolor{PrimaryRed}{\textbar}\quad Dr.\ sc., ETH Zurich\par
Germany / Switzerland / India \quad\textcolor{PrimaryRed}{\textbar}\quad \scholar
\par\smallskip
{\color{PrimaryRed}\hrule height 1pt}

\cvsection{Research Profile}
Atomistic simulations, metallic glasses, computational mechanics, finite element methods, and scientific computing. Research applies machine learning to materials discovery and connects materials science, photonics, and data-driven engineering.

\cvsection{Featured Research Publication}
\textbf{Joint research with Prof. Donna Strickland, 2018 Nobel Laureate in Physics}\par
S. P. Amirtharaj, Z. Xu, D. Strickland, B. Chowdhury, et al., including S. Swayamjyoti. \emph{Numerically optimized FROG results for the study of red-shifted spectra in multi-frequency Raman generation.} \href{https://arxiv.org/abs/2602.02463}{arXiv:2602.02463} (2026).\par
\textcolor{PrimaryRed}{\textbf{Role: Principal investigator (PI) and corresponding author.}}

\cvsection{Nobel Prize Engagement and Correspondence}
\begin{itemize}
\item \textbf{Nobel Prize Teacher Summit, Stockholm, Sweden} — Participated in 2024 and 2026. Received a personal invitation to the 2027 summit in Stockholm.
\item \textbf{IDEA League Mechanical Engineering Symposium, Aachen, Germany} — Received a letter from Prof. Joachim Frank, 2017 Nobel Laureate in Chemistry, for the symposium.
\end{itemize}

\cvsection{Current Position}
\textbf{Visiting Researcher}, Institute of General Mechanics, RWTH Aachen University, Germany. Research in computational mechanics, materials modelling, machine learning for engineering systems, scientific computing, and multiscale and data-driven simulation.

\cvsection{Academic Degrees}
\entry{Doctor of Science (Dr. sc.) in Materials Science}{2012–2016}
ETH Zurich, Switzerland. Joint research at the Laboratory of Metal Physics and Technology, ETH Zurich, and the Condensed Matter Theory Group, Paul Scherrer Institute.\par
\href{https://www.research-collection.ethz.ch/bitstream/handle/20.500.11850/169215/Doctoral_Thesis_of_Soumya_Swayamjyoti.pdf?sequence=1}{Doctoral thesis}: \emph{Atomistic simulations to study metallic glasses: a microscopic investigation of local structural excitations.}

\entry{M.Sc. in Computational Mechanics}{2009–2012}
Computational Mechanics of Materials and Structures, University of Stuttgart, Germany. One semester at UPC Barcelona, Spain.\par
\href{https://drive.google.com/file/d/1ZtLccqa9tsaxn4f3kyzWDGOhHQj0W_WQ/view}{Master's thesis}: \emph{Finite element implementation of orbital-free density functional theory.}

\entry{B.Tech. in Civil Engineering}{2004–2008}
NIT Karnataka, Surathkal, India. Thesis: \emph{Analysis of plates by Navier method.}

\newpage
\cvsection{Previous Academic and Professional Experience}
\begin{itemize}
\item \textbf{Scientific Researcher}, School of Minerals, Metallurgical and Materials Engineering,\newline IIT Bhubaneswar, India (2019–2025).
\item \textbf{Research Engineer}, Siemens Technology Center, Bangalore, India (2018).
\item \textbf{Software Engineer and AI Developer}, Hyscaler Consulting Group, Bhubaneswar, India (2018).
\item \textbf{Postdoctoral Research Associate}, School of Engineering, Brown University, USA (2017).
\item \textbf{Doctoral Researcher}, ETH Zurich and Paul Scherrer Institute, Switzerland (2012–2016).
\item \textbf{Graduate Student Assistant}, University of Stuttgart, Germany (2010–2012).
\item \textbf{Graduate Student Assistant}, CIMNE Barcelona, Spain (2009).
\item \textbf{Software Engineer}, Intergraph Process, Power and Marine, Hyderabad, India (2009).
\item \textbf{Software Engineer}, subsidiary of Exeter Group, Bangalore, India (2008).
\end{itemize}

\cvsection{Selected Publications}
See \scholar{} for the complete publication record and citation metrics.
\begin{enumerate}
\item S. Swayamjyoti, J. F. Löffler, and P. M. Derlet. \emph{Local structural excitations in model glass systems under applied load.} \href{https://doi.org/10.1103/PhysRevB.93.144202}{Physical Review B 93, 144202} (2016).
\item S. Swayamjyoti, J. F. Löffler, and P. M. Derlet. \emph{Local structural excitations in model glasses.} \href{https://doi.org/10.1103/PhysRevB.89.224201}{Physical Review B 89, 224201} (2014).
\item R. Kishore, S. Swayamjyoti, S. Das, A. K. Gogineni, Z. Nussinov, D. Solenov, and K. K. Sahu. \emph{Visual Machine Learning: Insight Through Eigenvectors, Chladni Patterns, and Community Detection in 2D Particular Structures.} International Journal for Multiscale Computational Engineering 20(3) (2022).
\item S. Das, H. Pegu, K. K. Sahu, A. K. Nayak, S. Ramakrishna, D. Datta, and S. Swayamjyoti. \emph{Machine Learning in Materials Modeling—Fundamentals and the Opportunities in 2D Materials.} Elsevier (2020). Principal investigator and corresponding author.
\item A. M. Menon, T. R. Kopparthi, P. Omprakash, H. Verma, A. Haldar, S. Swayamjyoti, K. K. Sahu, and C. Featherston. \emph{Deep Learning-Based Optimization of Piezoelectric Vibration Energy Harvesters.} AIAA SCITECH Forum (2022).
\item A. K. Gogineni, S. Swayamjyoti, D. Sahoo, K. K. Sahu, and R. Kishore. \emph{Multi-class classification of vulnerabilities in smart contracts using AWD-LSTM, with pre-trained encoder inspired from natural language processing.} IOP SciNotes 1(3), 035002 (2020).
\end{enumerate}

\cvsection{Research Funding and Grants}
\begin{itemize}
\item Swiss National Science Foundation doctoral research support, ETH Zurich and Paul Scherrer Institute (2012–2016), approximately CHF 200,000+.
\item Erasmus Mundus full scholarship, funded by the European Commission (2009–2012), approximately EUR 42,000.
\item Research project funding at IIT Bhubaneswar from India's Ministry of Human Resource Development and Ministry of Steel.
\end{itemize}

\newpage
\cvsection{Awards and Honours}
\begin{itemize}
\item Postdoctoral Fellowship, Brown University, USA.
\item Swiss National Science Foundation-funded doctoral researcher, ETH Zurich.
\item Erasmus Mundus Full Scholarship recipient.
\item Selected for the funded Entrepreneur First residency programme.
\end{itemize}

\cvsection{Teaching and Mentoring}
\begin{itemize}
\item Mentored doctoral, master's, and undergraduate students in machine learning, materials modelling, molecular simulations, and computational mechanics.
\item Co-instructed materials science and engineering courses and \emph{Introduction to Machine Learning for Engineers and Materials Scientists} for more than 100 undergraduates.
\end{itemize}

\cvsection{Conference Presentations and Workshops}
\begin{itemize}
\item \emph{Atomic-scale structural excitations in metallic glasses under zero and applied load}, IIT Madras, India (2019).
\item \emph{Microscopic structural excitations in model metallic glass systems}, Junior Euromat, Lausanne, Switzerland (2014).
\item CECAM Workshop, ETH Zurich, Switzerland (2013).
\item CCMX Summer School on Multi-scale Modelling of Materials, Switzerland (2012).
\item DUNE Workshop, Heidelberg, Germany (2011).
\end{itemize}

\cvsection{Invited Talks}
IIT Kharagpur, India (2019) \textperiodcentered{} Paul Scherrer Institute, Switzerland (2012) \textperiodcentered{} Max Planck Institute for Intelligent Systems, Germany (2011).

\cvsection{Leadership and International Academic Engagement}
Initiated the international collaboration with Prof. Donna Strickland's research group and led interdisciplinary collaborations involving researchers in Germany, Switzerland, India, Singapore, Canada, and the USA. Participated in Erasmus Mundus Association conferences and coordinated the venue for its General Assembly in Barcelona. Featured in EMANATE alumni success stories.

\cvsection{Research Areas and Technical Skills}
\textbf{Research:} Metallic glasses, molecular dynamics, atomistic simulations, computational mechanics, finite elements, numerical optimisation, machine learning for materials, AI for scientific discovery, and photonics.\par
\textbf{Computing:} Python, Fortran, C/C++, MATLAB, T-SQL, HTML/CSS/JavaScript, ASP.NET, Linux, high-performance computing, and deep learning.\par
\textbf{Industry:} Software engineering, AI systems, scientific computing, and enterprise technologies.

\cvsection{Languages and Career Breaks}
English (C2, fluent), Hindi and Oriya (native), German (A2). No career breaks.

\end{document}
